\chapter{Exercise Feedback}
\label{ch:exercise-feedback}
\label{app:exercise-feedback}

This appendix follows the exercise order in the reference.  Every heading gives
the exact source label and links back to the exercise.  ``Solution'' marks a
complete answer at the level requested.  ``Solution sketch'' gives the decisive
construction or invariant and an explicit criterion by which a completed proof
can be checked.

\section{Language levels, formal objects, and proof writing}

\paragraph{\hyperref[ex:levels-01]{\texttt{ex:levels-01}}.}
\textit{Solution.}  Relative to classical propositional logic, \(P\) and
\(P\land Q\) are object-language formulas, while \(\varphi\), the phrase ``is
valid,'' and \(\entails\) belong to the metalanguage.  A left-parenthesis token
used in the fully parenthesized formula \((P\land Q)\) is object-language
punctuation; the phrase ``the left parenthesis'' that mentions that token is
metalinguistic.

\paragraph{\hyperref[ex:levels-02]{\texttt{ex:levels-02}}.}
\textit{Solution.}  Put \(v(P)=\False\) and \(v(Q)=\True\).  Then
\(v(P\to Q)=\True\) and \(v(Q)=\True\), while \(v(P)=\False\).  Thus this
single row satisfies both premises and falsifies the conclusion.

\paragraph{\hyperref[ex:strings-01]{\texttt{ex:strings-01}}.}
\textit{Solution.}  Let \(B\) be the least set of strings such that
\(\epsilon\in B\), \(w\in B\) implies \((w)\in B\), and \(u,v\in B\) implies
\(uv\in B\).  Write \(L(w),R(w)\) for the two parenthesis counts.  Structural
induction gives \(L(\epsilon)=R(\epsilon)=0\); if \(L(w)=R(w)\), then both
counts increase by one in \((w)\); and if the equality holds for \(u,v\), then
\(L(uv)=L(u)+L(v)=R(u)+R(v)=R(uv)\).

\paragraph{\hyperref[ex:strings-02]{\texttt{ex:strings-02}}.}
\textit{Solution.}  Give every initial derivation height \(0\), and define the
height of a derivation ending with immediate subderivations
\(D_1,\ldots,D_k\) to be \(1+\max_i h(D_i)\).  Each immediate subderivation
has smaller height.  Repeating that strict inequality down the finite
constructor path shows that every proper subderivation has strictly smaller
height.

\paragraph{\hyperref[ex:writing-foundations-01]{\texttt{ex:writing-foundations-01}}.}
\textit{Solution.}  A correct completeness statement is: for every set
\(\Gamma\) of first-order sentences and sentence \(\sigma\), if
\(\Gamma\entails\sigma\) over all nonempty first-order structures, then
\(\Gamma\derives_{\ND_{\FOL}}\sigma\).  A false correction is: for every
arithmetic sentence \(\sigma\), if \(\mathcal N\entails\sigma\), then
\(\PA\derives\sigma\).  The first fixes a calculus and semantic class; the
second fails because true arithmetic is not the theorem set of \(\PA\).

\paragraph{\hyperref[ex:writing-foundations-02]{\texttt{ex:writing-foundations-02}}.}
\textit{Solution.}  The original leaves ``theory,'' ``true,'' the intended
structures, and the proof calculus unspecified.  Deductive completeness is:
for every sentence \(\sigma\), either \(T\derives\sigma\) or
\(T\derives\neg\sigma\).  The separate semantic property is that all models of
\(T\) agree on every sentence.  For a satisfiable \(T\), soundness and
completeness of the fixed first-order calculus make these two properties
equivalent.

\paragraph{\hyperref[ex:set-prerequisites-01]{\texttt{ex:set-prerequisites-01}}.}
\textit{Solution.}  Let \(i:A\to\N\) be an injection.  For each
\(b\in f[A]\), define
\[
  j(b)\defeq\min\set{i(a):a\in A\text{ and }f(a)=b}.
\]
The set minimized over is a nonempty subset of \(\N\).  If
\(j(b)=j(b')\), injectivity of \(i\) gives the same selected preimage, so
\(b=b'\).  Thus \(j:f[A]\to\N\) is an injection.  This also covers
\(A=\varnothing\), when \(j\) is the empty function.

\paragraph{\hyperref[ex:set-prerequisites-02]{\texttt{ex:set-prerequisites-02}}.}
\textit{Solution.}  Define \([a]_E+[b]_E=[a+b]_E\).  If
\(aE a'\) and \(bE b'\), then \(5\mid(a-a')\) and \(5\mid(b-b')\), so
\(5\mid((a+b)-(a'+b'))\).  Therefore \([a+b]_E=[a'+b']_E\), which proves
independence from both representatives.

\paragraph{\hyperref[ex:set-prerequisites-03]{\texttt{ex:set-prerequisites-03}}.}
\textit{Solution.}  For a family \((A_n)_{n\in\N}\), the usual proof first
selects, for every \(n\), one enumeration or surjection
\(e_n:\N\to A_n\) (with a separate convention for empty members), and only
then maps \((n,k)\) to \(e_n(k)\).  The simultaneous selection of the sequence
\((e_n)_n\) is the simultaneous selection step; pairing \(\N^2\) with \(\N\) is not.  If the
enumerations are part of the given data, that step is unnecessary.

\section{Classical propositional logic}

\paragraph{\hyperref[ex:prop-metatheory-01]{\texttt{ex:prop-metatheory-01}}.}
\textit{Solution.}  The full parse is
\(((\neg P\land Q)\to(R\lor S))\).  Its subformulas are
\(P,Q,R,S,\neg P,(\neg P\land Q),(R\lor S)\), and the whole implication.
Every binary node shown is fixed by the declared precedence, and implication
associativity is not needed because there is only one implication token.

\paragraph{\hyperref[ex:prop-metatheory-02]{\texttt{ex:prop-metatheory-02}}.}
\textit{Solution.}  Let \(l(\varphi)\) count syntax-tree leaves---occurrences
of propositional variables or \(\bottom\)---and let \(b(\varphi)\) count
binary-connective occurrences.  For a variable or \(\bottom\),
\(l=1=0+1=b+1\).  A negation has the same leaves and binary nodes as its sole
subformula, so the induction hypothesis is preserved.  If
\(\varphi\circ\psi\) is formed with a binary connective, then
\[
l(\varphi\circ\psi)=l(\varphi)+l(\psi)
=(b(\varphi)+1)+(b(\psi)+1)
=b(\varphi\circ\psi)+1.
\]
These are all constructors, so structural induction proves the invariant.

\paragraph{\hyperref[ex:prop-metatheory-03]{\texttt{ex:prop-metatheory-03}}.}
\textit{Solution.}  Set \(P=\False,Q=\True\).  Both \(P\to Q\) and \(Q\) are
true, but \(P\) is false, so semantic consequence fails.  This valuation
inspects truth values only; nonexistence of a derivation follows separately
from soundness of the specified calculus.

\paragraph{\hyperref[ex:prop-metatheory-04]{\texttt{ex:prop-metatheory-04}}.}
\textit{Solution.}  Let \(P\) mean that the process terminates, \(Q\)
that the input is finite, and \(R\) that the search is bounded.  Taking
``or'' as inclusive before applying the stated restriction, the phrase
following ``only if'' requires exactly one of \(Q\) and \(R\) whenever
\(P\) holds.  A faithful translation is
\[
  P\to\bigl((Q\lor R)\land\neg(Q\land R)\bigr).
\]
The stronger formula
\((P\to(Q\lor R))\land\neg(Q\land R)\) would forbid \(Q\land R\) even when
the process does not terminate; that global restriction is not stated by the
sentence.

\paragraph{\hyperref[ex:prop-metatheory-05]{\texttt{ex:prop-metatheory-05}}.}
\textit{Solution.}  A conjunctive normal form is
\((\neg P\lor Q)\land(\neg Q\lor R)\).  A disjunctive normal form is
\((\neg P\land\neg Q)\lor(Q\land R)\).  In the four rows with
\(Q=\False\), the formula is true iff \(P=\False\); in the four rows
with \(Q=\True\), it is true iff \(R=\True\).  Those two cases are
exactly the two displayed
disjuncts.

\paragraph{\hyperref[ex:prop-metatheory-06]{\texttt{ex:prop-metatheory-06}}.}
\textit{Solution.}  Assume \(P\to R\), then \(Q\to R\), then \(P\lor Q\).
For the left case assume \(P\) and infer \(R\); for the right case assume
\(Q\) and infer \(R\).  Disjunction elimination discharges the two case
assumptions.  Three successive applications of \(\to I\) discharge
\(P\lor Q\), \(Q\to R\), and \(P\to R\), in that order, yielding the required
nested implication.

\paragraph{\hyperref[ex:prop-metatheory-07]{\texttt{ex:prop-metatheory-07}}.}
\textit{Solution.}  Assume \(P\land\neg P\).  Conjunction elimination gives
\(P\) and \(\neg P\); negation elimination gives \(\bottom\); negation
introduction discharges the assumption and yields \(\neg(P\land\neg P)\).
The decisive final rule is \(\neg I\), which is available in both displayed
classical and intuitionistic natural deduction; no reductio rule is used.

\paragraph{\hyperref[ex:prop-metatheory-08]{\texttt{ex:prop-metatheory-08}}.}
\textit{Solution sketch.}  Apply \(\to R\) to reduce the target to
\(P\to Q,Q\to R,P\Rightarrow R\).  Apply \(\to L\) to \(Q\to R\); its right
premise is the identity sequent \(R\Rightarrow R\).  Its left premise derives
\(Q\) by applying \(\to L\) to \(P\to Q\), with identities
\(P\Rightarrow P\) and \(Q\Rightarrow Q\).  Context weakening supplies the
displayed companions.  A checked tree contains no cut, and every branch ends in
identity.

\paragraph{\hyperref[ex:prop-metatheory-09]{\texttt{ex:prop-metatheory-09}}.}
\textit{Solution sketch.}  Reflexivity, symmetry, and transitivity follow from
equality of truth values under every valuation.  For congruence, assume
corresponding inputs have equal truth values; each primitive connective is a
fixed truth function, so applying it to equal input tuples gives equal outputs.
Check the unary negation case and each binary truth table separately; these are
all primitive constructors.

\paragraph{\hyperref[ex:prop-metatheory-10]{\texttt{ex:prop-metatheory-10}}.}
\textit{Solution sketch.}  Choose the underivable formula \(P\).  Then
\(\set{\neg P}\) is consistent; extend it to a maximal consistent set
\(\Delta\).  Define \(v(Q)=\True\) iff \(Q\in\Delta\).  The truth
lemma, proved by induction on formulas, gives
\(v(\varphi)=\True\) iff \(\varphi\in\Delta\).  Since \(\neg P\in\Delta\),
maximal consistency gives \(P\notin\Delta\), hence \(v(P)=\False\).  A complete
answer must include the truth-lemma cases for every primitive connective.

\paragraph{\hyperref[ex:prop-metatheory-11]{\texttt{ex:prop-metatheory-11}}.}
\textit{Solution.}  The clause set is
\(\set{\set{P,Q},\set{\neg P,R},\set{\neg Q,R},\set{\neg R}}\).
Resolve the second and fourth clauses to obtain \(\set{\neg P}\), and the
third and fourth to obtain \(\set{\neg Q}\).  Resolve \(\set{P,Q}\) with
\(\set{\neg P}\) to get \(\set{Q}\), then with \(\set{\neg Q}\) to obtain
the empty clause.

\paragraph{\hyperref[ex:prop-metatheory-12]{\texttt{ex:prop-metatheory-12}}.}
\textit{Solution.}  A tautological clause is true under every valuation, so
removing it changes neither the satisfying valuations nor satisfiability.  For
the second claim take \(S=\set{\set{P},\set{Q}}\).  The valuation
\(P=Q=\True\) satisfies the clause \(\set{P}\) that is deleted, but the reduced
set \(\set{\set{Q}}\) is also satisfied by \(P=\False,Q=\True\), whereas
\(S\) is not.  Thus satisfaction by the current valuation does not establish
equivalence under all valuations.

\section{Classical first-order logic}

\paragraph{\hyperref[ex:fol-deduction-01]{\texttt{ex:fol-deduction-01}}.}
\textit{Solution.}  The free-variable set is \(\set{x,y,z}\).  The occurrence
of \(x\) in \(P(x,z)\) is free; the occurrences of \(x\) in \(Q(x,y)\) and
\(R(x,y,z)\) are bound by \(\forall x\).  The occurrence of \(y\) in
\(Q(x,y)\) is free because it lies outside the scope of \(\exists y\); the
occurrence of \(y\) in \(R(x,y,z)\) is bound by that existential quantifier.
Both occurrences of \(z\) are free.

\paragraph{\hyperref[ex:fol-deduction-02]{\texttt{ex:fol-deduction-02}}.}
\textit{Solution.}  In the first formula, rename the bound \(y\) to fresh
\(z\) before substitution, obtaining
\((\exists yR(x,y))[y/x]=\exists zR(y,z)\).  In the second, rename bound
\(x\) to fresh \(z\), because the replacing term contains \(x\); the result is
\((\forall x(P(x)\to Q(y)))[f(x)/y]
=\forall z(P(z)\to Q(f(x)))\).  In each result the formerly free variables in
the replacing term remain free.

\paragraph{\hyperref[ex:fol-deduction-03]{\texttt{ex:fol-deduction-03}}.}
\textit{Solution sketch.}  Prove simultaneously: term values agree whenever
two assignments agree on the term's variables, and formula truth agrees
whenever they agree on the formula's free variables.  The variable case is the
assignment hypothesis; the binary-function case applies both term induction
hypotheses; the binary-relation case applies term coincidence to both
arguments.  Boolean cases use formula induction.  In a quantifier case, update
both assignments at the bound variable by the same arbitrary element, then use
the induction hypothesis.  The proof for an arbitrary signature differs only
by replacing two arguments with the declared finite arity of the active symbol.

\paragraph{\hyperref[ex:fol-deduction-04]{\texttt{ex:fol-deduction-04}}.}
\textit{Solution.}  Choose \(y\) absent from both premises and the conclusion.
Under the assumption \(P(y)\), to be discharged by \(\exists E\),
universal elimination gives
\(P(y)\to Q(y)\); implication elimination gives \(Q(y)\); existential
introduction gives \(\exists xQ(x)\).  Apply \(\exists E\) to the premise
\(\exists xP(x)\) and discharge \(P(y)\).  The eigenvariable \(y\) is absent
from the conclusion, every other open assumption, and the existential premise.

\paragraph{\hyperref[ex:fol-deduction-05]{\texttt{ex:fol-deduction-05}}.}
\textit{Solution.}  From \(x=y\) and \(y=z\), apply equality elimination to
the latter equality with the formula \(\varphi(w)\defeq x=w\).  Its required
minor premise is \(\varphi(y)\), namely \(x=y\), and its conclusion is
\(\varphi(z)\), namely \(x=z\).  Thus \(x=y,y=z\derives x=z\).

\paragraph{\hyperref[ex:fol-deduction-06]{\texttt{ex:fol-deduction-06}}.}
\textit{Solution.}  Use one domain containing people and papers, unary
predicates \(\mathsf{Rev},\mathsf{Paper}\), and binary relations
\(\mathsf{Read},\mathsf{Cites}\).  The translations are
\[
\begin{aligned}
&\forall x\bigl(\mathsf{Rev}(x)\to\exists y(\mathsf{Paper}(y)\land
 \mathsf{Read}(x,y)\land
 \forall z((\mathsf{Paper}(z)\land\mathsf{Read}(x,z))\to z=y))\bigr),\\
&\exists y\bigl(\mathsf{Paper}(y)\land
 \forall x(\mathsf{Rev}(x)\to\mathsf{Read}(x,y))\bigr),\\
&\forall x\bigl(\mathsf{Rev}(x)\to\exists y(\mathsf{Paper}(y)\land
 \mathsf{Read}(x,y)\land\neg\mathsf{Cites}(y,y))\bigr).
\end{aligned}
\]
The last formula makes the nonvacuous reading explicit: every reviewer read at
least one paper that did not cite itself.

\paragraph{\hyperref[ex:fol-deduction-07]{\texttt{ex:fol-deduction-07}}.}
\textit{Solution.}  If a \(\Lang_{\mathrm{grp}}\)-structure satisfies the three
displayed sentences, the interpretation of \(\cdot\) is associative,
\(e\) is a two-sided identity, and \(i(a)\) is a two-sided inverse of each
\(a\); hence the reduct is a group.  Conversely, interpret \(\cdot,e,i\) in
any group by its multiplication, identity, and inverse operation.  The group
laws verify each universal sentence under every assignment.  Equality is
logical identity in both directions, so the model class is exactly the class of
groups in that signature.

\paragraph{\hyperref[ex:fol-deduction-08]{\texttt{ex:fol-deduction-08}}.}
\textit{Solution.}  In the signature with one binary relation \(E\), take
\[
T_{\mathrm{graph}}=\set{
\forall x\,\neg E(x,x),\quad
\forall x\forall y(E(x,y)\leftrightarrow E(y,x))}.
\]
Its models are precisely nonempty simple undirected graphs.  The first axiom
forbids loops, and logical identity needs no additional equivalence axioms.

\paragraph{\hyperref[ex:fol-deduction-09]{\texttt{ex:fol-deduction-09}}.}
\textit{Solution.}  Begin with the identity sequent \(P(y)\Rightarrow P(y)\).
Apply \(\forall L\) to obtain \(\forall xP(x)\Rightarrow P(y)\), then
\(\exists R\) to obtain
\(\forall xP(x)\Rightarrow\exists xP(x)\).  No cut occurs.  The semantic
soundness of choosing the instance denoted by \(y\) uses that assignments into
the domain exist, hence that the domain is nonempty.

\paragraph{\hyperref[ex:fol-cardinality-normalforms-01]{\texttt{ex:fol-cardinality-normalforms-01}}.}
\textit{Solution sketch.}  For \(\mathcal M\subseteq\mathcal N\), prove by
structural induction that formulas with parameters from \(M\) have the same
truth value in both structures.  Atomic and Boolean cases use the substructure
hypothesis and induction.  For \(\exists x\varphi\), the forward direction
keeps a witness in \(M\); the reverse direction uses the Tarski--Vaught witness
condition to replace an \(N\)-witness by one in \(M\), then applies induction.
For \(\forall x\varphi\), one direction restricts the quantified domain, while
the other argues that a counterexample in \(N\) would make
\(\exists x\neg\varphi\) true and hence have an \(M\)-witness.  A complete
proof must state both quantifier directions without invoking duality as a
substitute for these arguments.

\paragraph{\hyperref[ex:fol-cardinality-normalforms-02]{\texttt{ex:fol-cardinality-normalforms-02}}.}
\textit{Solution sketch.}  Let \(A_0\) be \(A\) together with \(m_0\)
and the interpretations of the constant symbols.  Given countable
\(A_n\), let \(A_{n+1}\) add the witness value \(F_\varphi(\bar a)\) for
every formula \(\exists y\,\varphi(y,\bar x)\) and tuple \(\bar a\) from
\(A_n\); the witness functions for the formulas
\(\exists y\,(y=f(\bar x))\) put the value of every signature function on
every tuple from \(A_n\) into \(A_{n+1}\).  These prescriptions explicitly
give \(A_0,A_1,A_2\), and countability follows from the countable signature.
If a finite tuple lies in \(\bigcup_nA_n\), it lies in some common stage; the
value of every signature function on that tuple lies in the next stage.  This
is the required closure check.

\paragraph{\hyperref[ex:fol-cardinality-normalforms-03]{\texttt{ex:fol-cardinality-normalforms-03}}.}
\textit{Solution.}  Add constants \(c_0,c_1,\ldots\) and all inequalities
\(c_i\ne c_j\) for \(i\ne j\) to \(T\).  Every finite subset has a model:
choose a sufficiently large finite model of the mentioned part of \(T\) and
interpret the finitely many constants distinctly.  Compactness gives a model
of the whole expanded theory, whose reduct is an infinite model of \(T\).  No
map from any particular finite model was constructed, so no embedding claim
follows.

\paragraph{\hyperref[ex:fol-cardinality-normalforms-04]{\texttt{ex:fol-cardinality-normalforms-04}}.}
\textit{Solution.}  Suppose an empty-signature sentence \(\sigma\) held exactly
in finite nonempty structures.  Add, for each \(n\geq1\), a sentence saying
that at least \(n\) distinct elements exist.  Every finite subset together with
\(\sigma\) has a sufficiently large finite model.  Compactness therefore gives
an infinite model of \(\sigma\), contradicting the assumed characterization.

\paragraph{\hyperref[ex:fol-cardinality-normalforms-05]{\texttt{ex:fol-cardinality-normalforms-05}}.}
\textit{Solution.}  Successively applying quantifier negation, implication
elimination, De Morgan's law, and vacuous-variable movement gives
\[
\begin{aligned}
&\neg\forall x(P(x)\to\exists y(Q(y)\land R(x,y)))\\
&\equiv\exists x\bigl(P(x)\land\neg\exists y(Q(y)\land R(x,y))\bigr)\\
&\equiv\exists x\forall y
  \bigl(P(x)\land(\neg Q(y)\lor\neg R(x,y))\bigr).
\end{aligned}
\]
The variables are distinct, and moving \(\forall y\) across the conjunction is
licensed because \(y\notin\FV(P(x))\).  The final formula is rectified and in
prenex normal form.

\paragraph{\hyperref[ex:fol-cardinality-normalforms-06]{\texttt{ex:fol-cardinality-normalforms-06}}.}
\textit{Solution.}  Replace \(u\) by a new constant \(c\); replace \(y\) by a
new unary function \(f(x)\); and replace \(w\) by a new binary function
\(g(x,z)\).  The universal closure is
\[
\forall x\forall z\,R(c,x,f(x),z,g(x,z)).
\]
The arities \(0,1,2\) record exactly the preceding universal variables on which
the corresponding existential witness may depend.

\paragraph{\hyperref[ex:fol-cardinality-normalforms-07]{\texttt{ex:fol-cardinality-normalforms-07}}.}
\textit{Solution.}  Let the original sentence be
\(\forall x\exists yR(x,y)\), and take the two-element structure with
\(R=\set{(0,0),(1,1)}\).  Its reduct satisfies the sentence.  Expanding by the
identity function makes \(\forall xR(x,f(x))\) true, while expanding by the
constant-zero function makes it false at \(1\).  Thus Skolemization is a claim
about existence of a suitable expansion and its reduct, not logical
equivalence of formulas from different signatures.

\paragraph{\hyperref[ex:fol-cardinality-normalforms-08]{\texttt{ex:fol-cardinality-normalforms-08}}.}
\textit{Solution.}  The first pair produces \(y=f(x)\) and
\(x=g(y)\), hence \(x=g(f(x))\); the occurs check proves there is no unifier.
The second pair has most general unifier \(\set{x\mapsto y}\) (equivalently,
with variables reversed, \(\set{y\mapsto x}\)).  The third produces
\(x=y\) and \(x=f(y)\), hence \(y=f(y)\); the occurs check again proves there
is no finite-term unifier.

\paragraph{\hyperref[ex:fol-cardinality-normalforms-09]{\texttt{ex:fol-cardinality-normalforms-09}}.}
\textit{Solution.}  After Skolemizing the existential premise with new constant
\(c\), the clauses are
\(\set{\neg P(x),Q(x)}\), \(\set{P(c)}\), and
\(\set{\neg Q(z)}\), with parent variables renamed apart.  Resolve the first
two using \(\set{x\mapsto c}\) to obtain \(\set{Q(c)}\).  Resolve that clause
with the third using \(\set{z\mapsto c}\) to obtain the empty clause.

\paragraph{\hyperref[ex:fol-cardinality-normalforms-10]{\texttt{ex:fol-cardinality-normalforms-10}}.}
\textit{Solution.}  A time bound truncates the search, so the result reports
only that no refutation was found among the inferences explored within that
resource limit.  Refutational completeness guarantees eventual discovery for
an unsatisfiable set only under an unbounded fair search.  It does not convert
finite unsuccessful search into a model or a satisfiability certificate.

\section{Model theory}

\paragraph{\hyperref[ex:maps-elementarity-01]{\texttt{ex:maps-elementarity-01}};
alias \hyperref[ex:diagram-test]{\texttt{ex:diagram-test}}.}
\textit{Solution.}  Interpret each new constant \(c_a\) in \(\mathcal N\) by
the same element \(a\in M\subseteq N\).  Substructure preservation of
constants and functions, and agreement of relations on tuples from \(M\), make
every true or negated atomic \(\Lang(M)\)-sentence from \(\mathcal M\) true in
the expansion of \(\mathcal N\).  The expansion satisfies the elementary
diagram iff, for every \(\Lang\)-formula \(\varphi(\bar x)\) and
\(\bar a\in M\),
\(\mathcal M\entails\varphi(\bar a)\) iff
\(\mathcal N\entails\varphi(\bar a)\), which is precisely
\(\mathcal M\preccurlyeq\mathcal N\).

\paragraph{\hyperref[ex:maps-elementarity-02]{\texttt{ex:maps-elementarity-02}}.}
\textit{Solution.}  Every formula in the language expanded by constants for
\(A\) becomes a formula of the original language with the corresponding
members of \(A\) as parameters.  Since \(\mathcal M\preccurlyeq\mathcal N\),
the two structures agree on that parameter instance.  Translating back gives
elementarity of the two constant expansions.

\paragraph{\hyperref[ex:maps-elementarity-03]{\texttt{ex:maps-elementarity-03}}.}\mbox{}\par
\noindent
\textit{Solution.}  Work in the equality-only signature.  Let
\(M_n=\set{0,1,\ldots,n}\), with inclusions as substructures.  Their union is
\(\N\), but \(M_0\) is not elementary in the union: with parameter \(0\), the
formula \(\exists x\,(x\ne0)\) is false in \(M_0\) and true in \(\N\).

\paragraph{\hyperref[ex:types-ultraproducts-01]{\texttt{ex:types-ultraproducts-01}};
alias \hyperref[ex:type-space-basis]{\texttt{ex:type-space-basis}}.}
\textit{Solution.}  A complete type contains both \(\varphi\) and \(\psi\)
iff it contains \(\varphi\land\psi\), giving the intersection
identity; it does not contain \(\varphi\) iff completeness puts \(\neg\varphi\)
in it, giving the complement identity.  For the equality-only theory of a
two-element structure, \(S_1(T)\) has one point, isolated by \([x=x]\).
The space \(S_2(T)\) has the equality type and inequality type, whose singleton
basic opens are \([x_1=x_2]\) and \([x_1\ne x_2]\), respectively.

\paragraph{\hyperref[ex:types-ultraproducts-02]{\texttt{ex:types-ultraproducts-02}};
alias \hyperref[ex:homogeneity-back-step]{\texttt{ex:homogeneity-back-step}}.}
\textit{Solution.}  The inverse \(h^{-1}:h[A]\to M\) is elementary.  Apply
\(\kappa\)-homogeneity to extend it to an elementary map
\(g:h[A]\cup\set{d}\to M\).  Elementary maps are injective, so
\(g^{-1}\) is elementary on its range; it extends \(h\), because
\(g(h(a))=a\), and its range contains \(d\), because
\(g^{-1}(g(d))=d\).  The relevant domains have size below \(\kappa\).

\paragraph{\hyperref[ex:types-ultraproducts-03]{\texttt{ex:types-ultraproducts-03}};
alias \hyperref[ex:omit-constant-type]{\texttt{ex:omit-constant-type}}.}
\textit{Solution.}  Suppose that \(\theta(x)\) were principal for
\(p(x)=\set{x\ne c_n:n\in\N}\).  Choose a model \(\mathcal M\models T\) and
\(a\in M\) with \(\mathcal M\models\theta(a)\).  If
\(a=c_k^{\mathcal M}\) for some \(k\), then this same witness shows that
\(T\not\entails\forall x(\theta(x)\to x\ne c_k)\), contrary to principality.
Otherwise \(a\ne c_n^{\mathcal M}\) for every \(n\).  The formula
\(\theta\) mentions only finitely many constants, so choose \(c_m\) that it
does not mention and form \(\mathcal M'\) by changing only the interpretation
of \(c_m\) to \(a\).  The interpretations of all constants remain pairwise
distinct, because \(a\) differed from every old constant interpretation;
hence \(\mathcal M'\models T\).  The truth of \(\theta(a)\) is unchanged,
but \(\mathcal M'\not\models a\ne c_m\), again contradicting principality.
Thus \(p\) is nonprincipal.  In the displayed model every element equals some
\(c_n\), so no element realizes \(p\).

\paragraph{\hyperref[ex:types-ultraproducts-04]{\texttt{ex:types-ultraproducts-04}}.}
\textit{Solution.}  Define
\(F([a]_{\mathcal U})=a(i_0)\).  Equality modulo an ultrafilter principal at
\(i_0\) means exactly equality at \(i_0\), so \(F\) is well-defined and
injective.  Every \(b\in M_{i_0}\) is the image of a choice function with
value \(b\) at \(i_0\), so \(F\) is surjective.  Functions and relations are
interpreted coordinatewise, hence \(F\) is an isomorphism.

\paragraph{\hyperref[ex:types-ultraproducts-05]{\texttt{ex:types-ultraproducts-05}}.}
\textit{Solution.}  In the ultrapower over \(\N\) by a nonprincipal
ultrafilter, let \(a=[i\mapsto i]_{\mathcal U}\).  For each standard \(n\),
the set \(\set{i:n<i}\) is cofinite and therefore belongs to
\(\mathcal U\).  By \L{}o\'s's theorem, \(d(n)<a\).  Thus \(a\) is greater
than every diagonal image while remaining one element of the ultrapower.

\paragraph{\hyperref[ex:qe-categoricity-01]{\texttt{ex:qe-categoricity-01}}.}
\textit{Solution sketch.}  Enumerate the domains of a countable dense order
without endpoints \(M\) and of \(\Q\).  Starting from the empty partial
isomorphism, alternately map the next unused element of \(M\) into the unique
open interval determined by its already mapped neighbors, and map the next
unused rational back into the corresponding interval of \(M\).  Density and
the absence of endpoints supply a fresh point in every step.  The union is
total, surjective, and order preserving and reflecting.  Check all four interval
cases: two bounds, lower bound only, upper bound only, and no bounds.

\paragraph{\hyperref[ex:qe-categoricity-02]{\texttt{ex:qe-categoricity-02}}.}
\textit{Solution.}  Given \(\varphi(\bar x)\), choose quantifier-free
\(\psi(\bar x)\) with
\(T\entails\forall\bar x(\varphi\leftrightarrow\psi)\).  Since both
structures model \(T\), each agrees with \(\psi\) on \(\varphi\).  A
substructure and its superstructure agree on quantifier-free formulas with
parameters from the substructure, so they agree on \(\varphi\).  If either
structure fails to model \(T\), the relative equivalence to \(\psi\) is not
available there, and the argument stops.

\paragraph{\hyperref[ex:qe-categoricity-03]{\texttt{ex:qe-categoricity-03}};
alias \hyperref[ex:back-and-forth-equivalence]{\texttt{ex:back-and-forth-equivalence}}.}
\textit{Solution sketch.}  Enumerate two countable models and maintain a finite
partial bijection preserving the equivalence relation.  To add an element
equivalent to an old domain element, choose a fresh member of the corresponding
infinite target class.  To add an element in a new class, choose a fresh target
class, possible because infinitely many classes remain.  Apply the symmetric
construction on alternate steps.  The union is a total isomorphism; checking
the two extension cases proves preservation and reflection of equivalence.

\section{Computability and arithmetical hierarchy}

\paragraph{\hyperref[ex:turing-machines-01]{\texttt{ex:turing-machines-01}}.}
\textit{Solution.}  Complement a decider by exchanging its two halting answers.
For finite union or intersection, run the finitely many deciders and combine
their answers by disjunction or conjunction.  For computably enumerable sets,
dovetail the recognizers: accept a union when one accepts, and accept an
intersection after every recognizer has accepted its common input.  No
complement construction is used, because a recognizer's failure to accept need
not halt.

\paragraph{\hyperref[ex:universal-reductions-01]{\texttt{ex:universal-reductions-01}}.}
\textit{Solution.}  Each property depends only on the computed partial
function, not on its program text, and is nontrivial.  The everywhere-undefined
function is not total while the zero function is; the former's range omits
\(0\) while the latter's contains \(0\).  Under the exercise's stated
nonempty-domain convention for a constant partial function, the zero function
has the third property and the everywhere-undefined function does not.  Rice's
theorem therefore makes all three index sets undecidable.

\paragraph{\hyperref[ex:decision-hierarchy-01]{\texttt{ex:decision-hierarchy-01}}.}
\textit{Solution sketch.}  With the fixed bounded arithmetic computation
formula \(\mathsf{Comp}(e,x,s)\), nonempty domain is defined by
\(\exists x\exists s\,\mathsf{Comp}(e,x,s)\), a \(\Sigma^0_1\) formula.  Given \(e\),
effectively produce a program that is undefined except that on input \(0\) it
simulates the normalized wrapper for program \(e\) on input \(e\), accepts
with a valid numerical output iff \(\varphi_e(e)\) is defined, and diverges
otherwise.
Its domain is nonempty iff \(e\in K\).  The uniform program-index map
is the required many-one reduction; writing it via the parameter theorem
completes the formal check.  Since \(K\) is
\(\boldsymbol{\Sigma}^0_1\)-complete under many-one reductions, this proves
the same completeness classification for the nonempty-domain index set.

\paragraph{\hyperref[ex:decision-hierarchy-02]{\texttt{ex:decision-hierarchy-02}}.}
\textit{Solution.}  Using the same bounded arithmetic formula
\(\mathsf{Comp}\), finite domain has
the immediate form
\(\exists b\forall x>b\forall s\,\neg\mathsf{Comp}(e,x,s)\), hence belongs to
\(\boldsymbol{\Sigma}^0_2\).  Infinite domain has form
\(\forall b\exists x>b\exists s\,\mathsf{Comp}(e,x,s)\), hence belongs to
\(\boldsymbol{\Pi}^0_2\).
Cofinite domain has form
\(\exists b\forall x>b\exists s\,\mathsf{Comp}(e,x,s)\), hence belongs to
\(\boldsymbol{\Sigma}^0_3\).
These are upper bounds only; no completeness claim follows without reductions.

\paragraph{\hyperref[ex:decision-hierarchy-03]{\texttt{ex:decision-hierarchy-03}}.}
\textit{Solution.}  Enumerate all finite derivations in a fixed effective,
sound, and complete first-order calculus and output their conclusions.  Every
output is valid by soundness.  Conversely, completeness is used exactly to
infer that every semantically valid sentence has some finite derivation and
therefore eventually appears.  The enumeration never ranges directly over
structures.

\section{Arithmetic and incompleteness}

\paragraph{\hyperref[ex:arithmetic-coding-01]{\texttt{ex:arithmetic-coding-01}};
alias \hyperref[ex:nonstandard-true-arithmetic]{\texttt{ex:nonstandard-true-arithmetic}}.}
\textit{Solution.}  Taking \(T=\Th(\mathcal N)\), the compactness construction
produces a countable \(\mathcal M\models T\) with an element unequal to every
numeral.  Since \(T\) contains exactly the sentences true in \(\mathcal N\),
\(\mathcal M\) and \(\mathcal N\) agree on every arithmetic sentence.  They
are not isomorphic, because every element of \(\mathcal N\) is denoted by a
numeral and the displayed element of \(\mathcal M\) is not.  Thus elementary
equivalence is strictly weaker than isomorphism.

\paragraph{\hyperref[ex:arithmetic-coding-02]{\texttt{ex:arithmetic-coding-02}}.}
\textit{Solution sketch.}  Fix primitive recursive coding of finite sequences
with length and component functions.  Let \(A(c,a,b)\) say: \(c\) has length
\(b+1\), its zeroth component is \(a\), and for every \(i<b\) its
\((i+1)\)-st component is the successor of its \(i\)-th component.  Bounded
quantification over components keeps \(A\) primitive recursive.  The last
component is \(a+b\); representability of \(A\), together with existence and
uniqueness of the coded recursion, yields a representing formula for addition.
A completed account must verify the sequence-code operations and the unique
last value.

\paragraph{\hyperref[ex:arithmetic-coding-03]{\texttt{ex:arithmetic-coding-03}}.}
\textit{Solution.}  The first assertion is a metalinguistic claim that a
\(T\)-derivation of \(\sigma\) exists.  The second is a semantic assertion in
the standard arithmetic structure about the numeral coding \(\sigma\).  The
third is again a metalinguistic derivability claim, now with an arithmetized
provability sentence as conclusion.  Correct representation of the proof
predicate connects an actual proof with standard truth of the second assertion;
if \(T\supseteq\mathsf Q\), numeralwise representability gives the first
derivability condition and connects an actual proof to the third.  The reverse
move from \(T\)-provability to standard truth requires the corresponding
soundness hypothesis.

\paragraph{\hyperref[ex:first-incompleteness-01]{\texttt{ex:first-incompleteness-01}}.}
\textit{Solution sketch.}  Effectiveness makes finite \(T\)-proof checking
primitive recursive and \(\mathsf{Prov}_T\) arithmetically expressible.
Representability turns a concrete proof into a \(T\)-proof of its proof-code
instance and verifies each standard nonproof code.  Consistency rules out
\(T\derives G_T\), since the fixed point would then yield both provability and
its negation.  \(\Sigma^0_1\)-soundness rules out
\(T\derives\neg G_T\), because \(\neg G_T\) is a false
\(\Sigma^0_1\) sentence once consistency shows \(G_T\) true in
\(\mathcal N\).  These four uses should appear at those exact transitions.

\paragraph{\hyperref[ex:first-incompleteness-02]{\texttt{ex:first-incompleteness-02}}.}
\textit{Solution.}  Consistency gives \(T\notderives G_T\).  Hence no standard
number codes a \(T\)-proof of \(G_T\), so
\(\mathcal N\entails\neg\mathsf{Prov}_T(\ulcorner G_T\urcorner)\).  The fixed
point is already provable in \(\mathsf Q\), which is true in \(\mathcal N\),
so \(\mathcal N\entails G_T\).  This semantic conclusion alone does not show
that \(T\cup\set{G_T}\) is consistent: that additionally requires
\(T\notderives\neg G_T\).  Under, for example, the stated
\(\Sigma^0_1\)-soundness hypothesis, that condition holds; then the one-sentence
extension is effective and consistent, so an incompleteness theorem supplies
another sentence it does not decide.  Thus adjoining \(G_T\) does not produce
all arithmetical truths under the strengthened hypothesis, while mere
consistency licenses no such claim about the extension.

\paragraph{\hyperref[ex:first-incompleteness-03]{\texttt{ex:first-incompleteness-03}}.}
\textit{Solution sketch.}  Once a proof with least standard code \(m\) is
assumed, every comparison with a code below \(m\) is a bounded finite check
whose true and false instances are numeralwise representable in \(\mathsf Q\);
\cref{lem:q-bounded-cases} turns the bounded existential over codes at most
\(m\) into a finite disjunction, and \cref{lem:q-dichotomy} splits an
arbitrary proof code against \(\bar m\) inside \(\mathsf Q\).
The Rosser sentence uses that bounded comparison to force an opposite proof
with no larger code, contradicting consistency.  Omega-consistency is not
needed because the argument never infers a universal standard claim from all
of its numeral instances; the least code converts the relevant search into a
single bounded verification.

\paragraph{\hyperref[ex:second-incompleteness-01]{\texttt{ex:second-incompleteness-01}}.}
\textit{Solution sketch.}  From the fixed point
\(G_T\leftrightarrow\neg\mathsf{Prov}_T(\ulcorner G_T\urcorner)\), internalize
the implication \(G_T\to\neg\mathsf{Prov}_T(\ulcorner G_T\urcorner)\) using
derivability conditions 1 and 2.  Condition 3 gives
\(\mathsf{Prov}_T(\ulcorner G_T\urcorner)\to
\mathsf{Prov}_T(\ulcorner
  \mathsf{Prov}_T(\ulcorner G_T\urcorner)\urcorner)\); internalized
modus ponens then yields
\(\mathsf{Prov}_T(\ulcorner G_T\urcorner)\to
\mathsf{Prov}_T(\ulcorner\bottom_A\urcorner)\).  Therefore
\(\mathsf{Con}(T)\) implies
\(\neg\mathsf{Prov}_T(\ulcorner G_T\urcorner)\), and the fixed point
gives \(\mathsf{Con}(T)\to G_T\).  If \(T\) proved its consistency, it would
prove \(G_T\), contradicting consistency.  A full derivation must display the
instances of all three conditions just identified.

\paragraph{\hyperref[ex:second-incompleteness-02]{\texttt{ex:second-incompleteness-02}}.}\mbox{}\par
\noindent
\textit{Solution.}  The assumed biconditional includes
\(T\derives\mathsf{Prov}_T(\ulcorner\varphi\urcorner)\to\varphi\).  This is
exactly the hypothesis of L\"ob's theorem, so
\(T\derives\varphi\).  The reverse half of the biconditional is not needed.

\paragraph{\hyperref[ex:second-incompleteness-03]{\texttt{ex:second-incompleteness-03}}.}
\textit{Solution.}  Suppose one formula \(\mathsf{Tr}(x)\) had all the stated
biconditionals provable in consistent \(T\supseteq\mathsf Q\).  By the diagonal
lemma choose \(L\) with
\(T\derives L\leftrightarrow
\neg\mathsf{Tr}(\overline{\ulcorner L\urcorner})\).  The
assumed instance for \(L\) gives
\(T\derives\mathsf{Tr}(\overline{\ulcorner L\urcorner})\leftrightarrow L\).
Propositional reasoning combines them to derive \(L\leftrightarrow\neg L\),
and hence a contradiction.  Consistency is used only at the final step to rule
out that derivation.

\paragraph{\hyperref[ex:second-incompleteness-04]{\texttt{ex:second-incompleteness-04}}.}
\textit{Solution.}  Three correct replacements are: (i) if \(T\) is a
consistent computably enumerable extension of \(\mathsf Q\), then its fixed
G\"odel sentence \(G_T\) is not provable in \(T\), and it is true in
\(\mathcal N\); (ii) if in addition the fixed provability predicate satisfies
the derivability conditions, then \(T\notderives\mathsf{Con}(T)\); and
(iii) \(\Th(\mathcal N)\) contains every true arithmetic sentence by
definition, is deductively complete, and is not computably enumerable.  These
claims concern different theories, predicates, and evidence.

\section{Set theory and forcing}

\paragraph{\hyperref[ex:zf-zfc-01]{\texttt{ex:zf-zfc-01}}.}
\textit{Solution sketch.}  Form \(a\cup b\) using Pairing and Union.  For
\(x\in a,y\in b\), the Kuratowski pair
\(\langle x,y\rangle=\set{\set{x},\set{x,y}}\) belongs to
\(\Pow(\Pow(a\cup b))\).  This double power set is the required bounding set.
Apply Separation to it with the formula
\(\exists x\in a\exists y\in b\,z=\langle x,y\rangle\).  Extensionality
identifies the separated set with \(a\times b\).  A complete derivation should
cite Pairing, Union, Power Set twice, Separation, and Extensionality at these
exact points.

\paragraph{\hyperref[ex:zf-zfc-02]{\texttt{ex:zf-zfc-02}}.}
\textit{Solution.}  By Pairing form the nonempty set \(c=\set{a,b}\).
Foundation supplies \(x\in c\) with \(x\cap c=\varnothing\).  If \(x=a\),
then \(b\in a\cap c\); if \(x=b\), then \(a\in b\cap c\).  Both cases
contradict the supplied disjointness, so no two-cycle \(a\in b\in a\) exists.

\paragraph{\hyperref[ex:cumulative-hierarchy-01]{\texttt{ex:cumulative-hierarchy-01}}.}
\textit{Solution.}  Fix \(\beta\) and argue directly from transitivity;
no transfinite induction is required.  If
\(\alpha=\beta\), inclusion is immediate.  If \(\alpha<\beta\), then
\(\alpha\in\beta\) by the ordinal order definition; since every ordinal
\(\beta\) is transitive, every member of \(\alpha\) is also a member of
\(\beta\), so \(\alpha\subseteq\beta\).  The use of transitivity is exactly the
step from \(u\in\alpha\in\beta\) to \(u\in\beta\).  The argument
invokes no induction hypothesis, which is why no transfinite induction is
required.

\paragraph{\hyperref[ex:cumulative-hierarchy-02]{\texttt{ex:cumulative-hierarchy-02}}.}
\textit{Solution.}  Put \(X=\set{\varnothing,\set\varnothing}\).  Since
\(\rank(\varnothing)=0\) and \(\rank(\set\varnothing)=1\),
\[
\rank(X)=\sup\set{0+1,1+1}=2.
\]
Also \(X\subseteq V_2\), so \(X\in V_3\), while
\(\set\varnothing\notin V_1\), so \(X\notin V_2\).  This confirms the rank
calculation from the hierarchy stages.

\paragraph{\hyperref[ex:cumulative-hierarchy-03]{\texttt{ex:cumulative-hierarchy-03}}.}
\textit{Solution.}  Let \(X=\set{a,b}\) with \(a\ne b\) and let
\(R=\varnothing\).  The relation is well-founded but not extensional, because
\(a\) and \(b\) have the same predecessor set.  The collapse recursion sends
both to \(\set{}=\varnothing\), so the collapse map is not injective.

\paragraph{\hyperref[ex:constructibility-forcing-01]{\texttt{ex:constructibility-forcing-01}}.}
\textit{Solution.}  For a finite partial function \(p\), if
\(n\in\dom(p)\) then \(p\in D_n\); otherwise
\(p\cup\set{(n,0)}\) extends \(p\) and belongs to \(D_n\).  Thus \(D_n\) is
dense.  If \((n,i)\in p\in G\) and \((n,j)\in q\in G\), directedness of the
filter gives \(r\in G\) extending both \(p,q\).  Since \(r\) is a function,
\(i=j\).  Hence \(\bigcup G\) is a function; meeting every \(D_n\) further
makes its domain all of \(\omega\).

\paragraph{\hyperref[ex:constructibility-forcing-02]{\texttt{ex:constructibility-forcing-02}}.}
\textit{Solution sketch.}  \(L\entails\mathsf{CH}\) is an external semantic
claim about the constructible universe, proved from the fine structure and
definable well-order of \(L\).  The formula
\(\ZFC\derives V=L\to\mathsf{CH}\) is a syntactic derivability claim obtained
by formalizing that result in \(\ZFC\).  The implication
\(\mathsf{Con}(\ZFC)\to\mathsf{Con}(\ZFC+\mathsf{CH})\) is a
metatheoretic relative-consistency claim: from a model or syntactic consistency
construction for \(\ZFC\), pass to its constructible inner model, which
satisfies \(V=L\) and hence \(\mathsf{CH}\).  A completed proof must state the
soundness/completeness or syntactic formalization used when moving between
model existence and consistency.

\section{Sequent calculus, normalization, and typed proof terms}

\paragraph{\hyperref[ex:lk-lj-01]{\texttt{ex:lk-lj-01}}.}
\textit{Solution sketch.}  Apply \(\land L\) to the antecedent, making
\(A\) and \(B\lor C\) available; then apply \(\lor L\) to \(B\lor C\).  In
the \(B\)-branch, identities for \(A,B\) and \(\land R\) yield
\(A\land B\), followed by the left \(\lor R\) injection.  In the
\(C\)-branch, the analogous steps yield \(A\land C\), followed by the right
injection.  The principal formulas, from bottom upward, are
\(A\land(B\lor C)\), \(B\lor C\), the concluding disjunction, and the
relevant conjunction.  The finished tree contains no cut and has one succedent at every
node.

\paragraph{\hyperref[ex:lk-lj-02]{\texttt{ex:lk-lj-02}}.}
\textit{Solution.}  Without the eigenvariable condition, identity
\(P(y)\Rightarrow P(y)\) would permit \(\forall R\) to infer
\(P(y)\Rightarrow\forall xP(x)\).  Take domain \(\set{0,1}\), interpret
\(P\) as \(\set{0}\), and assign \(y\) to \(0\).  The antecedent is true under
that assignment and the succedent is false.  Thus the relaxed rule derives an
invalid sequent.

\paragraph{\hyperref[ex:cut-elimination-01]{\texttt{ex:cut-elimination-01}}.}
\textit{Solution sketch.}  Suppose the left cut premise obtains
\(A\land B\) by \(\land R\) from derivations of \(A\) and \(B\), while the
right cut premise obtains its sequent by \(\land L\) from a derivation with
both \(A,B\) in the antecedent.  First cut the derivation of \(B\) against the
right premise to remove \(B\), retaining \(A\); then cut the derivation of
\(A\) against that result.  Context contractions or exchanges are applied as
licensed by the active calculus.  Both new cut formulas, \(A\) and \(B\), are
proper subformulas of \(A\land B\), so each has strictly smaller cut rank.

\paragraph{\hyperref[ex:cut-elimination-02]{\texttt{ex:cut-elimination-02}}.}
\textit{Solution.}  In propositional \(\LJ\), a cut-free backwards search uses
only finitely many subformulas of the endsequent; with the declared context
representation, only finitely many relevant sequents can arise, so exhaustive
search with repetition control decides derivability.  In first-order
\(\LJ\), quantifier rules introduce arbitrarily many term instances, and the
set of possible terms is generally infinite.  Literal subformula finiteness
therefore does not by itself bound proof search.

\paragraph{\hyperref[ex:normalization-curry-howard-01]{\texttt{ex:normalization-curry-howard-01}}.}
\textit{Solution.}  The term has type \(A\to(B\to A)\).  In the typing
calculus, the hypothesis \(y:B\) enters the context and is never consumed
by the \textsc{Var} rule; that unused context entry is the structural
weakening implicit in \(\lambda x{:}A.\lambda y{:}B.x\).  In the tree
natural-deduction calculus, a discharging application must discharge at
least one open occurrence, so the corresponding derivation of
\(A\to(B\to A)\) routes the assumption \(B\) through a conjunction
detour: assume \(A\) and \(B\), form \(B\land A\) by \(\land I\), recover
\(A\) by \(\land E_2\), discharge \(B\) to obtain \(B\to A\), and
discharge \(A\).  Both discharged occurrences are genuinely used by
\(\land I\).

\paragraph{\hyperref[ex:normalization-curry-howard-02]{\texttt{ex:normalization-curry-howard-02}}.}
\textit{Solution sketch.}  The required typing substitution lemma is: if
\(\Theta_{\mathrm{ty}},x{:}A\derives M:B\) and
\(\Theta_{\mathrm{ty}}\derives N:A\), then
\(\Theta_{\mathrm{ty}}\derives M[N/x]:B\), after capture-avoiding
renaming of bound variables.  Typing inversion on
\(\Theta_{\mathrm{ty}}\derives(\lambda x{:}A.M)N:B\) gives its two
premises; the lemma then types the reduct \(M[N/x]\) by \(B\).  Capture
avoidance is used in the lambda-abstraction case of the lemma when the
binder occurs free in \(N\).
Checking that case and the variable case completes the outer-redex proof.

\paragraph{\hyperref[ex:normalization-curry-howard-03]{\texttt{ex:normalization-curry-howard-03}}.}
\textit{Solution.}  If every reduction sequence from a term terminates, then in
particular at least one terminating sequence exists, so strong normalization
implies weak normalization.  For the converse failure, take a relation with
\(a\to a\) and \(a\to b\), where \(b\) has no successor.  The object \(a\) has
the terminating path \(a\to b\), but also the infinite path
\(a\to a\to a\to\cdots\).

\section{Second-order, typed, intuitionistic, and modal systems}

\paragraph{\hyperref[ex:second-order-logic-01]{\texttt{ex:second-order-logic-01}}.}
\textit{Solution.}  Let \(X\subseteq\N\) be arbitrary, not assumed definable,
and suppose \(0\in X\) and \(n\in X\Rightarrow n+1\in X\).  External induction
on \(n\) gives \(n\in X\) for every \(n\in\N\), hence \(X=\N\).  Since full
second-order semantics ranges its unary-relation variables over every subset of
\(\N\), this verifies the full induction axiom in the standard structure.

\paragraph{\hyperref[ex:second-order-logic-02]{\texttt{ex:second-order-logic-02}}.}
\textit{Solution.}  The categoricity proof forms the range of the recursively
defined map \(f:\N\to M\) and instantiates second-order induction with that
range to show it is all of \(M\).  In a Henkin structure, the unary-relation
domain may omit \(\operatorname{ran}(f)\).  The induction axiom then cannot be
instantiated with that subset, so the surjectivity step fails.

\paragraph{\hyperref[ex:second-order-logic-03]{\texttt{ex:second-order-logic-03}}.}
\textit{Solution sketch.}  Any finitary strongly complete calculus would make
full second-order consequence compact: from \(\Gamma\entails\varphi\), strong
completeness gives a finite derivation, which uses only a finite
\(\Gamma_0\subseteq\Gamma\), so \(\Gamma_0\entails\varphi\) by soundness.  The
displayed full-semantics counterexample to compactness contradicts this.  This
argument excludes finitary sound strong completeness; the stronger claim that
full second-order validities are not computably enumerable requires a separate
effectivity theorem.

\paragraph{\hyperref[ex:simple-type-theory-01]{\texttt{ex:simple-type-theory-01}}.}
\textit{Solution.}  Here
\(R:\iota\to\iota\to o\), \(x:\iota\),
\(Rx:\iota\to o\), and \(Rxx:o\).  Therefore
\(\lambda x{:}\iota.Rxx:\iota\to o\), and the whole term has type
\((\iota\to\iota\to o)\to(\iota\to o)\).

\paragraph{\hyperref[ex:simple-type-theory-02]{\texttt{ex:simple-type-theory-02}}.}
\textit{Solution.}  Under standard semantics,
\(\forall f{:}(\alpha\to\beta).A\) is true under assignment \(s\) iff
\(A\) is true under \(s[f\mapsto g]\) for every function
\(g:D_\alpha\to D_\beta\).  Under Henkin semantics the same clause quantifies
only over the designated domain \(D_{\alpha\to\beta}\), whose members need not
be all such functions.  The quantifier range is the only difference in the two
clauses.

\paragraph{\hyperref[ex:simple-type-theory-03]{\texttt{ex:simple-type-theory-03}}.}
\textit{Solution.}  A precise true statement is: the standard typed deductive
calculus is sound and complete for its declared class of Henkin structures.  A
precise false statement is: some finitary effective sound calculus is strongly
complete for standard full higher-order semantics.  Thus the bare word
``complete'' omits both the calculus and the semantic class and can conceal a
change from Henkin to full semantics.

\paragraph{\hyperref[ex:intuitionistic-logic-01]{\texttt{ex:intuitionistic-logic-01}}.}
\textit{Solution.}  Let \(w<v\), and let \(\mathcal M\) make \(A\) true only
at \(v\) and \(B\) true nowhere.  At both \(w\) and \(v\), \(A\to B\) fails
because \(\mathcal M,v\entails A\) and \(\mathcal M,v\notentails B\).
Hence \(\mathcal M,w\entails(A\to B)\to A\) vacuously, while
\(\mathcal M,w\notentails A\).  Peirce's formula therefore fails at \(w\).

\paragraph{\hyperref[ex:intuitionistic-logic-02]{\texttt{ex:intuitionistic-logic-02}}.}
\textit{Solution.}  If \(\mathcal M,w\entails\neg(A\lor B)\), then no
extension of \(w\) satisfies \(A\lor B\); consequently no extension satisfies
\(A\), and none satisfies \(B\).  Thus
\(\mathcal M,w\entails\neg A\land\neg B\).  Conversely, if
\(\mathcal M,w\entails\neg A\land\neg B\), an extension satisfying
\(A\lor B\) would satisfy one disjunct and contradict the corresponding
negation.  Hence the converse also holds intuitionistically.

\paragraph{\hyperref[ex:modal-temporal-epistemic-01]{\texttt{ex:modal-temporal-epistemic-01}}.}
\textit{Solution.}  Let \(W=\set{w,u}\) and
\(R=\set{(w,w),(w,u),(u,u)}\).  This frame is reflexive and is not Euclidean,
since \(wRw\) and \(wRu\) but not \(uRw\).  Let \(\mathcal M\) be the
structure on this frame that makes \(p\) true only at \(w\).
Then \(\mathcal M,w\entails\Diamond p\), but
\(\mathcal M,u\notentails\Diamond p\); because \(wRu\),
\(\mathcal M,w\notentails\Box\Diamond p\).

\paragraph{\hyperref[ex:modal-temporal-epistemic-02]{\texttt{ex:modal-temporal-epistemic-02}}.}
\textit{Solution.}  At positions \(0,1,2\), choose position \(3\) as the
witness for \(q\); \(p\) holds at every earlier position, so
\(p\mathbin{\mathsf U}q\) is true.  It is also true at position \(3\), with
the current position as witness and an empty intervening interval.  Its value
at position \(4\) is undetermined from the data: one must specify the truth of
\(p,q\) there and on the infinite continuation, because a past occurrence of
\(q\) cannot witness a future-until formula.

\paragraph{\hyperref[ex:modal-temporal-epistemic-03]{\texttt{ex:modal-temporal-epistemic-03}}.}
\textit{Solution.}  Take agents \(a,b\) and worlds \(w,u,v\).  Let the
\(a\)-equivalence classes be \(\set{w,u}\), \(\set v\), and the
\(b\)-classes be \(\set w\), \(\set{u,v}\).  Make \(p\) true at \(w,u\) and
false at \(v\).  At \(w\), both agents know \(p\), but the two-step path
\(w\sim_a u\sim_b v\) reaches a world where \(p\) is false.  Therefore
\(C_{\set{a,b}}p\) fails at \(w\).

\section{Many-valued, relevance, and linear logic}

\paragraph{\hyperref[ex:many-valued-paraconsistent-01]{\texttt{ex:many-valued-paraconsistent-01}}.}
\textit{Solution sketch.}  For \(\mathbf K_3\), the two roots are
\((p\land q){:}1,p{:}0\) and
\((p\land q){:}1,p{:}\tfrac12\).  Expanding the conjunction labelled \(1\)
forces \(p{:}1,q{:}1\), so each root closes on conflicting labels for \(p\).
For \(\mathsf{LP}\), test designated premise labels \(1,\tfrac12\) against
the sole nondesignated conclusion label \(0\).  Every branch of the
\(\min\)-expansion of a designated conjunction gives \(p\) a designated
value, so every branch again conflicts with \(p{:}0\).  Listing all
\(\min(a,b)=\tfrac12\) input pairs completes the LP tree.

\paragraph{\hyperref[ex:many-valued-paraconsistent-02]{\texttt{ex:many-valued-paraconsistent-02}}.}
\textit{Solution.}  Use \(p\lor\neg p\).  At values \(0,\tfrac12,1\), it has
values \(1,\tfrac12,1\), respectively.  All three are designated in
\(\mathsf{LP}\), whereas \(\tfrac12\) is not designated in
\(\mathbf K_3\).  Hence the formula is LP-valid but not K3-valid.

\paragraph{\hyperref[ex:relevance-linear-01]{\texttt{ex:relevance-linear-01}}.}
\textit{Solution.}  Write \(g\) for the tensor of \(\Gamma\).  The premise of
\(\multimap R\) is \(g\otimes a\leq b\), which is equivalent by residuation to
\(g\leq a\multimap b\), its conclusion.  For \(\multimap L\), suppose
\(g\leq a\) and \(d\otimes b\leq c\).  Since
\(a\otimes(a\multimap b)\leq b\), monotonicity gives
\(d\otimes g\otimes(a\multimap b)\leq d\otimes b\leq c\), exactly the
conclusion sequent's inequality.

\paragraph{\hyperref[ex:relevance-linear-02]{\texttt{ex:relevance-linear-02}}.}
\textit{Solution.}  Use identities \(A\Rightarrow A\) and
\(B\Rightarrow B\) as the two premises of \(\multimap L\).  It yields
\(A,A\multimap B\Rightarrow B\), and exchange built into multiset contexts
gives \(A\multimap B,A\Rightarrow B\).  Every antecedent occurrence is consumed
by the rule, so no weakening is required.

\paragraph{\hyperref[ex:relevance-linear-03]{\texttt{ex:relevance-linear-03}}.}
\textit{Solution.}  Sequence contexts build in none of exchange, weakening, or
contraction.  Multiset contexts identify permutations and thus build in
exchange, but retain multiplicity and do not build in weakening or contraction.
Set contexts ignore order and multiplicity, so exchange and contraction are
built in; weakening is still a substantive monotonicity principle because a
set does not identify \(\Gamma\) with \(\Gamma\cup\set A\).

\paragraph{\hyperref[ex:relevance-linear-04]{\texttt{ex:relevance-linear-04}}.}
\textit{Solution.}  The decisive clause says that \(A\to B\) at \(w\) tests
every pair \(u,v\) with \(R(w,u,v)\), requiring truth of \(A\) at \(u\) to
yield truth of \(B\) at the potentially different world \(v\).  In the
displayed structure, \(R(a,u,v)\), \(q\) holds at \(u\), and \(p\) fails at
\(v\), so \(q\to p\) fails at \(a\), which then defeats the outer
conditional at the normal world.  An ordinary binary Kripke clause evaluates
antecedent and consequent at one accessible world and therefore does not encode
this ternary split between the two tested worlds.

\section{Formal practice, automation, mechanization, and typesetting}

\paragraph{\hyperref[ex:formal-practice-01]{\texttt{ex:formal-practice-01}}.}
\textit{Solution.}  On a one-element domain \(\set{a}\), the premise makes
\(P(a)\lor Q(a)\) true, so one disjunct is true; that disjunct then holds of
every domain element, making one universal true.  For a two-element
countermodel, take domain \(\set{a,b}\), let \(P=\set{a}\) and
\(Q=\set{b}\).  Every element satisfies \(P\lor Q\), but neither predicate
holds of every element, so the conclusion is false.

\paragraph{\hyperref[ex:formal-practice-02]{\texttt{ex:formal-practice-02}}.}
\textit{Solution sketch.}  Induct structurally on the formula, after proving
the term-substitution lemma by induction on terms.  Strengthen the claim to all
structures, assignments, variables, and replacing terms at once:
\(\mathcal M,s\entails\varphi[t/x]\) iff
\(\mathcal M,s[x\mapsto\sem{t}^{\mathcal M}_s]\entails\varphi\).  Atomic and
Boolean cases use the corresponding induction hypotheses.  In
\(\forall y\psi\) and \(\exists y\psi\), if \(y=x\) substitution stops below
the binder; if \(y\ne x\) but \(y\in\operatorname{Var}(t)\), first rename
\(y\) fresh.  A
completed proof must show the two assignment updates commute after that
renaming.

\paragraph{\hyperref[ex:automated-reasoning-01]{\texttt{ex:automated-reasoning-01}}.}
\textit{Solution.}  The unit clause \(\set{\neg R}\) sets \(R=\False\).
Then \(\set{\neg P,R}\) and \(\set{\neg Q,R}\) reduce to units
\(\set{\neg P}\) and \(\set{\neg Q}\), setting \(P=Q=\False\).  The first
clause \(\set{P,Q}\) then becomes empty.  DPLL returns \(\UNSAT\).

\paragraph{\hyperref[ex:automated-reasoning-02]{\texttt{ex:automated-reasoning-02}}.}
\textit{Solution.}  Decomposing the two atoms gives
\(x\doteq g(y)\) and \(x\doteq y\).  Substituting the second equation into the
first gives \(y\doteq g(y)\).  Since \(y\) occurs properly in \(g(y)\), the
occurs check fails, and no finite-term unifier exists.

\paragraph{\hyperref[ex:automated-reasoning-03]{\texttt{ex:automated-reasoning-03}}.}
\textit{Solution.}  Before Skolemization, equivalence-preserving steps give
\[
\neg\forall x(P(x)\to\exists yR(x,y))
\equiv\exists x\bigl(P(x)\land\forall y\neg R(x,y)\bigr).
\]
Replace the existential witness by a new constant \(c\).  This step preserves
only equisatisfiability between the original signature and suitable expansions.
The remaining universal formula is equivalent in the expanded signature to the
universal closures of the clauses \(\set{P(c)}\) and
\(\set{\neg R(c,y)}\).  Thus only the Skolem step weakens equivalence to
equisatisfiability.

\paragraph{\hyperref[ex:mechanization-01]{\texttt{ex:mechanization-01}}.}
\textit{Solution.}  Suitable terms are
\[
\lambda x{:}A.x,\qquad
\lambda p{:}(A\times B).\pi_1p,\qquad
\lambda f{:}(A\to B).\lambda g{:}(B\to C).\lambda x{:}A.g(fx).
\]
They have types \(A\to A\), \((A\times B)\to A\), and
\((A\to B)\to(B\to C)\to(A\to C)\), respectively, by variable,
product-elimination, and function-elimination typing rules.

\paragraph{\hyperref[ex:mechanization-02]{\texttt{ex:mechanization-02}}.}
\textit{Solution sketch.}  In a deep embedding, let
\(\mathsf{Form}_{d}\) be an inductive datatype with constructor
\(\mathsf{And}:\mathsf{Form}_{d}\to\mathsf{Form}_{d}\to\mathsf{Form}_{d}\),
and define \(\mathsf{eval}_{d}(A,v)\) by recursion.  The required adequacy
theorem is
\[
  \mathsf{eval}_{d}(A,v)=\mathsf{true}
  \quad\Longleftrightarrow\quad
  v\entails_{\mathrm{CPL}} A
\]
for every object-language formula \(A\) and valuation \(v\); its conjunction
case uses
\(\mathsf{eval}_{d}(\mathsf{And}(A,B),v)=
 \mathsf{eval}_{d}(A,v)\land_{\mathrm{Bool}}\mathsf{eval}_{d}(B,v)\).

For a shallow embedding, fix for each valuation \(v\) a translation
\(\llbracket-\rrbracket^{s}_{v}\) from object-language formulas to host
propositions and stipulate
\[
  \llbracket A\land B\rrbracket^{s}_{v}
  \defeq
  \llbracket A\rrbracket^{s}_{v}\land_{\mathrm{host}}
  \llbracket B\rrbracket^{s}_{v}.
\]
Its required adequacy theorem is
\[
  \llbracket A\rrbracket^{s}_{v}
  \quad\Longleftrightarrow\quad
  v\entails_{\mathrm{CPL}}A
\]
for every \(A\) and \(v\).  The deep representation retains syntax for
induction, substitution, and quotation; the shallow representation supplies
those operations only if additional encoded syntax is introduced.

\paragraph{\hyperref[ex:latex-style-01]{\texttt{ex:latex-style-01}}.}
\textit{Solution.}  Replace ``Let \(f\) be the inclusion.  Clearly atomic
formulas are preserved.  Clearly the result follows by induction'' with:
``Let \(f:M\hookrightarrow N\) be the inclusion map of the
\(\Lang\)-substructure \(\mathcal M\subseteq\mathcal N\).  For every function
symbol \(g\), substructure closure gives
\(f(g^{\mathcal M}(\bar a))=g^{\mathcal N}(f\bar a)\); for every relation
symbol \(R\), the substructure definition gives
\(R^{\mathcal M}(\bar a)\) iff \(R^{\mathcal N}(f\bar a)\).  These are the
atomic induction cases.  The Boolean constructor cases follow from their truth
clauses and the induction hypotheses.''  The typed map and the two omitted
justifications are now explicit.

\paragraph{\hyperref[ex:latex-style-02]{\texttt{ex:latex-style-02}}.}
\textit{Solution.}  For a first-order resolution chapter, register
\(x,y,z\) as object variables, \(r,t\) as terms,
\(A,B\) as atoms, \(C,D\) as finite literal clauses, and
\(\theta_{\mathrm u},\rho_{\mathrm u},\eta_{\mathrm u}\) as finite-support
term substitutions.  Register \(\operatorname{ucl}_{\Lang}(C)\) as the formula
interpretation of a clause.  Do not also use \(C\) for a calculus or
\(\theta_{\mathrm u}\) for a formula in that chapter.  If \(\entails\) occurs,
its argument shapes must separately declare structure satisfaction and clause
set consequence, as in the registered-notation appendix.

\section{Keys for the qualifying problems}

The five labelled qualifying problems retain their existing solutions.  This
list makes their keys part of the same deterministic lookup system.

\paragraph{\hyperref[prob:substitution]{\texttt{prob:substitution}}.}
\textit{Solution.}  The solution immediately following the problem
identifies the term-substitution lemma, sets up the structural induction,
and displays the quantified assignment-commutation calculation; the
term-lemma induction and the atomic and connective cases are left to the
reader.

\paragraph{\hyperref[prob:compactness]{\texttt{prob:compactness}}.}
\textit{Solution.}  The complete solution immediately following the problem
builds models for finite subsets, applies compactness, and separates the
existential model claim from the rejected assertion about all orders.

\paragraph{\hyperref[prob:incompleteness]{\texttt{prob:incompleteness}}.}
\textit{Solution.}  The solution immediately following the problem states the
 computably enumerable axiom-set, arithmetic-strength, and consistency hypotheses and
distinguishes the original G\"odel conclusion from Rosser's refinement.

\paragraph{\hyperref[prob:modal-frame]{\texttt{prob:modal-frame}}.}
\textit{Solution.}  The complete solution immediately following the problem
proves the reflexive direction by the satisfaction clause and constructs a
valuation at a nonreflexive world for the converse.

\paragraph{\hyperref[prob:cut-subformula]{\texttt{prob:cut-subformula}}.}
\textit{Solution.}  The solution immediately following the problem traces the
subformula invariant through cut-free propositional rules, derives consistency
as a consequence, and states the required first-order qualification modulo term
substitution and eigenvariable renaming.

